\section{Where each wins, and why}
\label{sec:analysis}

\subsection{The mechanism behind each gap}

\begin{itemize}
  \item \textbf{Proof size --- recursion.} FRI recursively folds the
  codeword ($\log N$ halvings, a few hashes per round), giving a
  $\polylog$ proof. Brakedown opens $t\approx987$ full columns flat (no
  recursion on a random RAA code), giving $t\cdot\sqrt N$ with a large
  $t$-floor. \zinc{} bought a linear-time encoder; the price is the flat,
  large proof.
  \item \textbf{Verifier --- who reads the circuit.} \binius{}'s verifier
  evaluates the wiring multilinear in the clear, $O(N)$
  (\Cref{sec:shifts}). \zinc{}'s uniform AIR makes the structure an
  $O(\log N)$ formula, so its verifier pays only the $\sqrt N$ proximity
  test.
  \item \textbf{Memory / large $N$ --- reduction field.} \binius{}'s
  $\mathrm{GF}(2^{128})$ tensors plus a rate-$\tfrac12$ codeword spill
  cache super-linearly; \zinc{}'s base-field, bit-packed discharge streams
  $16$--$32\times$ fewer bytes and scales sub-linearly.
  \item \textbf{Prover --- a tie.} Both discharge the same $\AND$ content
  with cheap encoders; the word-level versus bit-sliced asymmetry roughly
  cancels (\Cref{sec:res-phases}).
\end{itemize}

\subsection{Could \binius{} take an AIR to make its verifier succinct?}
\label{sec:ana-air}

The natural counterfactual: since \binius{}'s verifier is $O(N)$ only
because of its arithmetization, could it adopt a uniform AIR on the SHA
path and recover a succinct verifier? The answer is \emph{yes for SHA in
isolation} --- the result is a binary-field STARK --- but with caveats
that explain why \binius{} does not.

\paragraph{What is actually $O(N)$.} Crucially, \binius{}'s
\emph{commitment} verification is already $\polylog$ (FRI). The $O(N)$
cost is entirely the structure/wiring check (the monster sweep,
$\sim$$99.7\%$ of verify at $2^{22}$). So a verifier fix is a
structure-check fix; the FRI commitment can stay. A uniform AIR collapses
the structure check two ways: cross-step dataflow becomes a fixed
next-row shift ($O(\log N)$), and within-word rotations become fixed
virtual columns ($O(\wbit)$, via the same shift indicator). The
combination ``binary commitment $+$ sumcheck $+$ uniform AIR $+$ FRI'' is
coherent --- it is a binary STARK --- and would be succinct on the
verifier and small on the proof.

\paragraph{Quantified estimate.} Decomposing the measured numbers (not a
built implementation): removing the monster sweep leaves the FRI residual,
so verify would be $\sim$$0.5$--$1$\,ms at $2^{20}$ rising only to
$\sim$$2$--$3$\,ms at $2^{24}$ (near-flat, $\polylog$) --- a $\sim$$25\times$
to $\sim$$1000\times$ win over current \binius{}, and faster than \zinc{}'s
$\sqrt N$ verify at scale. The prover stays a tie (the AIR removes the
shift/wiring prover but adds a wider padded trace and selector columns;
neither dominates --- and there is no prover speedup, consistent with the
prover and verifier being decoupled). The proof stays small (FRI is kept).

\paragraph{Why \binius{} does not.} (i) \emph{Generality}: \binius{} is a
general non-deterministic word-circuit framework; an AIR only covers
uniform computations, so the rewrite abandons the framework's reason for
being. (ii) The general verifier fix is therefore not an AIR but a
\emph{holographic / Spark} structure commitment \cite{spartan,lasso} ---
commit the wiring once, prove its evaluations via a sparse polynomial
commitment in $O(\log N)$, keeping the general circuit. (iii) An AIR
re-introduces a rigid power-of-two grid and selector columns for SHA's
irregular parts (the once-per-block feed-forward add, the
first-$16$-vs-last-$48$ message-schedule split, the padding/length
binding \binius{} proves in-circuit), and a wider committed trace than its
lean non-deterministic witness.

\begin{remark}[The same shifts, succinctly]
A uniform AIR makes the entry-wise shift succinct \emph{on the verifier}
regardless of word representation, because a uniform shift touches only
the constant intra-word subcube. But to get \zinc{}-clean shift handling
\emph{on the prover} as well (a monomial read-off rather than a shift
reduction), one needs the $\Ftwo[X]$ cell ring or a bit-decomposed layout
--- i.e. ``\binius{} done well on the prover'' converges toward \zinc{}'s
arithmetization (\Cref{sec:shifts-repr}).
\end{remark}

\subsection{Could a well-optimised \zinc{} win the prover outright?}
\label{sec:ana-prover}

\zinc{} has four structural prover advantages, all flowing from the
$\Ftwo[X]$ representation doing less algebraic work per gate: (i) the
entry-wise shift is a free $\psi$ read-off rather than a shift-reduction
sumcheck; (ii) a linear-time RAA encoder (no NTT); (iii) one commitment
serving both $\psia$ and $\psiz$; (iv) base-field discharge with a
bit-packed witness, giving the memory/scaling edge. Advantage (iv) is the
strongest and already wins large $N$.

The honest scoreboard, however, is a regime split rather than a clean
sweep:
\begin{itemize}
  \item At moderate $N$ on adequate RAM, \binius{} is $\sim$$1.05$--$1.13\times$
  faster --- its mature word-level prover and cheaper commit outweigh
  \zinc{}'s advantages, which are partly cancelled by the
  $\times\wbit$-vs-word-level asymmetry.
  \item At large $N$, \zinc{}'s lower per-compression floor
  ($\sim$$0.078$--$0.083$ vs $\sim$$0.090$\,ms) and sub-linear scaling pull
  it ahead; the crossover is around $2^{21}$--$2^{22}$.
\end{itemize}
There is one real blocker to stating ``\zinc{} wins the prover'' flatly:
the soundness asymmetry (\Cref{sec:cmp-soundness}). A fair, both-sound
comparison must add \zinc{}'s sound-adder cost, which is prover time
\zinc{} has not yet paid. \emph{If} that cost is roughly prover-neutral
(plausible --- the lookup adder removes $12$ Hadamards while adding one
grand product, but it is an open A/B), then a well-optimised, fully sound
\zinc{} edges \binius{} at scale and ties at moderate $N$ --- a scaling
win, not a universal one. The cleanest experiment is to ship the
soundness fix and the residual word-level optimisation and run both
provers fully sound, back-to-back, on a host with adequate RAM.
